All 1,564 maintained tests pass in 301.283 seconds; the 24 focused tests
pass in 0.734 seconds. The nine new regressions check parallel meridians,
vertex links, layered sources, odd/even one-sided families, connected normal
doubles, source/weight/base/count mutations, shared trace ownership, caps,
operation limits and cancellation. Binary JSON and a 16,385-bit multiplicity
pass independent replay without expanding components.

The 68.207-second source audit computes inventories on all 5,241 saved
cut-surface instances on 48 triangulations. Every core count, full normal
base vector, multiplicity and Euler value agrees with explicit finite chamber
recovery using the independent geometric pairing construction. The inventories
account for 7,838 midsection instances. This is agreement of complete
coordinate inventories, not just agreement of total component counts.
Representative bases on all 48 sources receive fresh Regina connectivity
and Euler checks. The earlier geometric pilot supplies the full finite
midsection/cut-signature context; no unique 3-manifold classification is
inferred from those signatures alone.

All native proofs replay with weight, inventory, chamber and orbit producers
disabled. Six large cut-orbit traces are then reused with orbit search
disabled; their inventories agree exactly with fresh queries. Nine large
binary cases include one-, four- and sixteen-tetrahedron meridian families
at $2^{128}$ or $2^{500}$ scale, the $2^{16384}$ one-tetrahedron case, and
large odd/even M\"obius stacks. The meridian families have one base vector
with multiplicity $m-1$. Even M\"obius stacks retain both the one-sided
vector and its connected annular double where present; odd stacks retain
the annular vector. The audit preserves 57 representative complete source
proofs.

For the large meridian families, maximum observed weight-run counts are
six, 24 and 96 for one, four and sixteen source tetrahedra respectively;
they do not grow between the tested $2^{128}$ and $2^{500}$ multiplicities.
These are finite diagnostic observations, not the proof of the polynomial
bound. That bound comes from the source weight construction, normal-type
accounting and the existing weighted interval operation.

Five paired timing rounds measure 80 completed inventory calls and sixteen
warmups, with identical-baseline A/A arms. Both methods return the same full
midsection vectors and multiplicities. The explicit reference expands the
finite chamber graph, computes its connected components and accumulates
normal-disc counts for each unmarked component. The native arm constructs
compressed indicators, runs weighted transport, independently replays the
complete source proof and serializes its result. All source construction,
validation, reconstruction and serialization are charged. This is an explicit
graph baseline, not an optimized published parallelity-bundle constructor.
\begin{center}
\begin{tabular}{rrrrrr}
\toprule
copies & explicit ms & native ms & ratio & old A/A & new A/A\\
\midrule
16 & 0.168 & 0.880 & 0.187 & 0.982 & 1.000\\
256 & 0.575 & 0.897 & 0.641 & 0.995 & 1.000\\
4096 & 8.060 & 0.928 & 8.647 & 1.011 & 0.989\\
65536 & 154.142 & 1.033 & 148.591 & 0.991 & 1.073\\
\bottomrule
\end{tabular}
\end{center}


At 65,536 copies the paired explicit/native ratio is 148.591, with median
154.142 versus 1.033 milliseconds. At 4,096 copies the ratio is 8.647.
At sixteen and 256 copies the ratios are 0.187 and 0.641: native proof and
weighted reconstruction lose on small inputs. Old A/A ratios range from
0.982 to 1.011 and new ones from 0.989 to 1.073. These remain host observations,
not universal speedup factors or full-recognition timings.

Runtime release \code{a10e8eb83}, all tests, exact source inventories,
search-free trace reuse checks and complete-query timings are retained.

The reproduction driver is \code{normal\_orbit\_research.prismatic\_inventory}.
Records are under \code{data/prismatic-inventory-*}. Publication review checks
664 current source/evidence hashes, the exact pilot input, every timing outcome
and the passing test records. It independently replays all 57 saved source
proofs with producers disabled. The complete recognizer is unchanged. The
next geometric task is parallelity inside the marked cores and its attachment
representation; that richer problem is not replaced by this whole-component
inventory.
